% STATUS: DRAFT
\chapter{Lorentzian signature and open problems}\label{ch:lorentzian_nc}
\skippable

\begin{physmotiv}
Physics is Lorentzian; spectral triples are Euclidean. This chapter lists what is known and not about moving to Lorentzian signature, and other open issues of the program.
\end{physmotiv}

\begin{itemize}
\item \textbf{Krein spaces.} A Lorentzian Dirac operator is not self-adjoint on a Hilbert space but is self-adjoint with respect to an indefinite inner product on a Krein space; several Lorentzian spectral triple proposals exist, with axioms differing from the Euclidean ones.
\item \textbf{Wick rotation.} In general no analytic continuation exists on a time-dependent background; the spectral action in its heat-kernel form is Euclidean. A spectral \emph{Lorentzian} action and its asymptotic expansion are not established in generality.
\item \textbf{Causal structure.} Distance formulas have Lorentzian analogues (Moretti, Franco--Eckstein, and others) involving a causal cone of states; no consensus.
\item \textbf{Quantization.} The spectral action is classical; coupling to quantum matter fields and gravitational quantization remains unclear.
\item \textbf{Unitarity and positivity} of the resulting field theory with higher-derivative terms from $a_4$.
\end{itemize}
(These remarks summarize the field's open directions and should be checked against surveys.)

\section*{Exercises}
\begin{exercise}
Why does the heat kernel $\Tr\ee^{-tD^2}$ fail to exist for a Lorentzian Dirac operator? (Hint: $D^2$ is not bounded below.)
\end{exercise}
